% modules/open-targets/A4-variational-inference.tex
% -----------------------------------------------------------------------------
% Deep-dive companion to the A4 entry of the Part I agenda (modules/08-open-targets):
% transport constants for variational inference -- turning reverse-KL/ELBO
% optimization into certified geometric (W_2) and posterior-mean error bars.
% Compiles inside main.tex; standalone it shows cross-module \ref as "??".
% -----------------------------------------------------------------------------
\documentclass[../../main.tex]{subfiles}
\begin{document}
\section{A4 --- Transport constants for variational inference}
\label{sec:a4}

This is the deep-dive companion to the A4 paragraph of the agenda (Section~\ref{sec:agenda}).
Variational Bayes (VB) reports the minimizer of a \emph{reverse} Kullback--Leibler divergence,
$q^\star\in\arg\min_{q\in\mathcal Q}\KL(q\,\|\,\pi)$, but the user cares about a geometric error
--- the displacement $\Wtwo(q^\star,\pi)$, or the posterior-mean error
$\norm{\E_{q^\star}\theta-\E_\pi\theta}$. The orientation is exactly right: reverse KL is the
quantity a Talagrand $T_2$ inequality (Definition~\ref{def:tci}) converts into transport,
\begin{equation}\label{eq:a4-bridge}
  \pi\ \text{satisfies}\ T_2(\CTCI)
  \quad\Longrightarrow\quad
  \Wtwo^2(q,\pi)\ \le\ 2\,\CTCI\,\KL(q\,\|\,\pi)
  \qquad\text{for all } q\ll\pi,
\end{equation}
so a transport constant is the natural ``conversion rate'' from an ELBO/KL error into a certified
posterior error. The A4 slogan --- \emph{small reverse KL} $+$ \emph{good transport constant}
$\Rightarrow$ \emph{small geometric error} --- is therefore correct. The whole content is hidden
in the word ``good'': the constant must be of \emph{posterior} scale, not prior scale, and a
global $T_2$ constant is a sufficient but often far too crude certificate. The purpose of this
section is to make precise the three increasingly relevant constants A4 proposes (global,
family-restricted, first-moment), to record where the naive $T_2$ route is wasteful or simply
infinite, and to state the certification theorem one should aim for.

\noindent\textbf{Reviewed status.}
The unrestricted mean dual, its localized refinement, the three distinct local-VI expansions,
the exact Gaussian calibration, the symmetrization lemma, and the global logistic-family rigidity
statement have passed independent review. The multimodal lower-witness calculations remain an
explicit analytic draft. The final posterior-scale GLM certificate remains an open question.

\subsection{Setup: the variational objective and the transport bridge}
\label{subsec:a4-setup}

Let $\pi(\dd\theta)\propto e^{-U(\theta)}\dd\theta$ be the posterior and $\mathcal Q$ a tractable
variational family (mean-field, Gaussian, low-rank Gaussian, structured mean-field, a
normalizing-flow family). VB solves
\[
  q^\star\in\arg\min_{q\in\mathcal Q}\KL(q\,\|\,\pi),
  \qquad
  \mathrm{ELBO}(q)\ =\ \log Z-\KL(q\,\|\,\pi),\quad Z=\int e^{-U},
\]
optimizing the ELBO in place of the intractable $\KL$ \cite{BleiKucukelbirMcAuliffe2017}. The
existing statistical theory of VB is rich but answers a \emph{different} question: it certifies
risk bounds, posterior-contraction rates, consistency, and asymptotic normality of variational
posteriors \cite{WangBlei2019,ZhangGao2020,YangPatiBhattacharya2020,RaySzabo2020}, not a
computable inequality of the form $\Wtwo(q,\pi)\le(\text{number})$ or
$\norm{\E_q\theta-\E_\pi\theta}\le(\text{number})$. The adjacent optimal-transport VI literature
also answers a different question --- how a VI \emph{algorithm} behaves in Wasserstein geometry,
or what happens when one minimizes a Wasserstein objective \emph{instead} of KL
\cite{AmbrogioniEtAl2018,YiLiu2023,LambertEtAl2022,ArneseLacker2024,JiangChewiPooladian2023}.
A4 asks the post-hoc certification question: \emph{given} the ordinary KL-based objective, when
does small KL certify small $\Wtwo$ or small mean error?

The figure of merit is a hierarchy of worst-case conversion factors. The global $T_2$ constant
(Definition~\ref{def:tci}) is the unrestricted one,
\begin{equation}\label{eq:a4-ctci}
  \CTCI(\pi)\ =\ \sup_{\substack{q\ll\pi\\ \KL(q\|\pi)>0}}\frac{\Wtwo^2(q,\pi)}{2\,\KL(q\,\|\,\pi)} .
\end{equation}
Because VB never ranges over all $q\ll\pi$, the operative quantity is the
\emph{family-restricted} constant
\begin{equation}\label{eq:a4-cq}
  C_{\mathcal Q}(\pi)\ :=\ \sup_{\substack{q\in\mathcal Q\\ \KL(q\|\pi)>0}}
  \frac{\Wtwo^2(q,\pi)}{2\,\KL(q\,\|\,\pi)}
  \ \le\ \CTCI(\pi),
\end{equation}
the inequality being a restriction of the supremum to $\mathcal Q$. Since actual VB solutions
live near the posterior mass, the sharpest object is the \emph{localized} restricted constant
\begin{equation}\label{eq:a4-cqr}
  C_{\mathcal Q,r}(\pi)\ :=\ \sup_{\substack{q\in\mathcal Q\\ 0<\KL(q\|\pi)\le r}}
  \frac{\Wtwo^2(q,\pi)}{2\,\KL(q\,\|\,\pi)} ,
\end{equation}
which only sees perturbations within ELBO reach of the optimum. Write
\begin{equation}\label{eq:a4-gap}
  \delta_{\mathcal Q}:=\inf_{q\in\mathcal Q}\KL(q\,\|\,\pi).
\end{equation}
The sublevel is empty for $r<\delta_{\mathcal Q}$, so local asymptotics should be parameterized by
$r=\delta_{\mathcal Q}+\rho$, not by an unqualified ``small $r$.'' The certificate these constants
deliver is
\begin{equation}\label{eq:a4-cert}
  \Wtwo^2(q^\star,\pi)\ \le\ 2\,C_{\mathcal Q,\,\KL(q^\star\|\pi)}(\pi)\,\KL(q^\star\,\|\,\pi).
\end{equation}

\begin{remark}[A constant alone is not an error bar]
\label{rem:a4-elbo}
Equation~\eqref{eq:a4-cert} is computable only with a numerical upper bound on
$\KL(q^\star\,\|\,\pi)$, and the ELBO supplies $\KL=\log Z-\mathrm{ELBO}$, which requires
$\log Z$. A complete VB error-bar theory therefore needs \emph{two} ingredients --- a transport
constant \emph{and} a computable KL certificate (a rigorous upper bound on $\log Z$, obtained by
a separately justified model-specific method). The A4 constants are the first ingredient; the second is
orthogonal and should be stated as a separate hypothesis, not folded into $C_{\mathcal Q}$.
\end{remark}

\subsection{Front 1: the family-restricted constant for log-concave and GLM posteriors}
\label{subsec:a4-restricted}

For a strongly log-concave posterior the global constant is already of posterior scale: if
$\Hess U\succeq mI$ then Bakry--\'Emery and Otto--Villani
(Theorems~\ref{thm:bakry-emery}, \ref{thm:otto-villani}) give $\CTCI(\pi)\le1/m$, and for a
Gaussian $N(\mu,\Sigma)$ this is sharp, $\CTCI=\lmax(\Sigma)$
(Family~\ref{fam:gaussian-covariance}). For a regular parametric posterior with
$\Sigma\approx(nI(\theta_0))^{-1}$ this is the right scale,
$\CTCI\asymp\tfrac1n\lmax(I(\theta_0)^{-1})$ --- exactly the scale VB error bars want, and the
$n\to\infty$ statement of A2 (Conjecture~\ref{conj:a2}, Equation~\eqref{eq:a2-target}).

The Gaussian-prior GLM posterior of Section~\ref{sec:tier3} is where the global constant remains
at the wrong scale. Theorem~\ref{thm:glm-fi} certifies $\CTCI\le\lmax(\Sigmaz)$, but the floor
$m=\lmax(\Sigmaz)^{-1}$ is the \emph{prior} curvature: the likelihood Hessian
$\sum_i\ell''(x_i^\top\theta)\,x_ix_i^\top$ degenerates in the tails
(Proposition~\ref{prop:flat-prior-nonintegrable}), so the global $T_2$ constant is exactly
prior-scale for binary logistic regression (Proposition~\ref{prop:a4-logistic-global}), even
though the posterior concentrates at the likelihood scale
$(\Sigmaz^{-1}+X^\top\bar W X)^{-1}$. This is precisely the gap that makes the \emph{localized}
constant the useful object: a mean-field Gaussian $q$ near the VB optimizer lives in the
high-curvature bulk, so $C_{\mathcal Q,\delta_{\mathcal Q}+\rho}(\pi)$ can be of posterior scale
even when both $C_{\mathcal Q}$ and $\CTCI$ are prior-scale. The A4 target is therefore the $\Wtwo$-analogue of the data-informed
A1 question (Question~\ref{q:glm-data-informed}): a sharp restricted/localized $C$ for the GLM
posteriors of Section~\ref{sec:tier3}, of order
$\lmax[(\Sigmaz^{-1}+X^\top\bar W X)^{-1}]$ in the bulk.

\begin{warning}[Restricted constants are not automatically finite --- heavy tails]
\label{warn:a4-heavytail}
A4 is right that $C_{\mathcal Q}$ can be far smaller than $\CTCI$, but it is \emph{not}
automatically finite, even for a Gaussian location family. Take the one-dimensional target
$\pi(\dd x)\propto e^{-\abs{x}^p}\dd x$ with $1\le p<2$ (so $\pi$ has lighter-than-exponential
but heavier-than-Gaussian tails for $p$ near $1$, and exactly exponential at $p=1$), and the
Gaussian location family $q_m=N(m,\sigma^2)\in\mathcal Q$. As $\abs m\to\infty$,
\[
  \begin{aligned}
    \Wtwo^2(q_m,\pi)&\asymp m^2,
    &\KL(q_m\,\|\,\pi)&=\E_{q_m}\abs{x}^p+O(1)\asymp\abs m^p,\\
    \text{so}\qquad
    \frac{\Wtwo^2(q_m,\pi)}{\KL(q_m\,\|\,\pi)}&\asymp\abs m^{2-p}\longrightarrow\infty.
  \end{aligned}
\]
Hence $C_{\mathcal Q}(\pi)=\infty$ for $p<2$ already at the mean-field Gaussian level (for
polynomial-tailed $\pi$ the failure is stronger, $\KL$ growing only logarithmically in the
location while $\Wtwo^2$ grows quadratically). A finite localized constant thus needs a coercive
family on the KL sublevel, quadratic-exponential tail control of $\pi$, or a weaker/modified
transport cost (Subsection~\ref{subsec:a4-modified}). A KL cutoff alone is not sufficient over a
rich family: putting mass $\eps_R\asymp r/R^p$ on $[R,R+1]$ keeps KL below $r$ while making the
second moment of order $rR^{2-p}\to\infty$. Thus even an unrestricted entropy ball can have
infinite $W_2$ radius.
This is the $\Wtwo$ shadow of the LSI/$T_2$ obstruction of Tier~4
(Theorem~\ref{thm:heavy-tail-no-lsi}, Family~\ref{fam:heavy-tail-classification}).
\end{warning}

\subsection{Front 2: the linear-mean entropy constant is the right primary object}
\label{subsec:a4-mean}

VB users report $\E_q\theta$, not $q$ itself, so the operative error is
$\norm{\E_q\theta-\E_\pi\theta}$. Throughout this subsection take
$\mathcal Q\subset\mathcal P_1(\R^d)$ and restrict to finite-entropy competitors. The $T_2$ route
bounds the error,
\begin{equation}\label{eq:a4-t2-to-mean}
  \norm{\E_q\theta-\E_\pi\theta}\ \le\ W_1(q,\pi)\ \le\ \Wtwo(q,\pi)\ \le\ \sqrt{2\,\CTCI\,\KL(q\,\|\,\pi)},
\end{equation}
but wastefully: $\Wtwo$ penalizes variance, tail, and covariance mismatch, whereas a mean error
only tests \emph{linear} observables. The matching constant is
\begin{equation}\label{eq:a4-cmean}
  C_{\mathrm{mean},\mathcal Q}(\pi)\ :=\ \sup_{\substack{q\in\mathcal Q\\ \KL(q\|\pi)>0}}
  \frac{\norm{\E_q\theta-\E_\pi\theta}^2}{2\,\KL(q\,\|\,\pi)}
  \ =\ \sup_{\norm u=1}\ \sup_{\substack{q\in\mathcal Q\\ \KL(q\|\pi)>0}}
  \frac{\big(u^\top(\E_q\theta-\E_\pi\theta)\big)^2}{2\,\KL(q\,\|\,\pi)} .
\end{equation}
Its localized version $C_{\mathrm{mean},\mathcal Q,r}$ restricts the same supremum to
$0<\KL(q\|\pi)\le r$.

\begin{theorem}[Exact unrestricted posterior-mean dual]
\label{thm:a4-mean-dual}
Assume that $\pi$ has a finite first moment. With extended values allowed,
\begin{equation}\label{eq:a4-mean-dual}
 C_{\mathrm{mean},\mathrm{all}}(\pi)
 =\sup_{\norm u=1}\sup_{t\ne0}
 \frac{2}{t^2}\log\E_\pi e^{t u^\top(\theta-\E_\pi\theta)},
\end{equation}
where $\log(+\infty)=+\infty$.
\end{theorem}

\begin{proof}
For unrestricted $q$, the mean constant is governed by the linear functionals of $\pi$ through
the Donsker--Varadhan / Gibbs variational principle:
$\E_q g\le\KL(q\,\|\,\pi)+\log\E_\pi e^g$. Taking
$g=t\,u^\top(\theta-\E_\pi\theta)$ and assuming the directional Laplace bound
$\log\E_\pi\exp(tu^\top(\theta-\E_\pi\theta))\le K(u)t^2/2$ gives
\[
  t\,u^\top(\E_q\theta-\E_\pi\theta)-\tfrac12K(u)t^2\ \le\ \KL(q\,\|\,\pi)\quad\forall t,
\]
and optimization over $t$ yields the mean inequality with constant $K(u)$. Conversely, inserting
the unrestricted mean inequality with constant $C$ into the Gibbs variational formula and
optimizing over the entropy gives the same Laplace bound with $K(u)=C$. Taking the two suprema
proves \eqref{eq:a4-mean-dual}, including the extended-valued case.
\end{proof}

This theorem is independently agent-certified in
\path{solutions/eq-a4-mean-dual.tex}.
For a restricted family only
$C_{\mathrm{mean},\mathcal Q}\le C_{\mathrm{mean},\mathrm{all}}$ is automatic. This
linear-observable inequality is weaker than a genuine $T_1$ inequality, which controls all
Lipschitz observables, and much weaker than $T_2$; it is tied only to the sub-Gaussian Laplace
bounds of the reported linear functionals \cite{BolleyVillani2005}.

The same entropy projection gives an exact localized formula, rather than merely the global
sub-Gaussian envelope. For $\norm u=1$ put
\[
 X_u=u^\top(\theta-\E_\pi\theta),\qquad
 \psi_u(t)=\log\E_\pi e^{tX_u},\qquad
 r_u(t)=t\psi_u'(t)-\psi_u(t).
\]

\begin{proposition}[Localized unrestricted mean tilts]
\label{prop:a4-mean-local}
Assume that, for every unit $u$ for which $X_u$ is nonconstant, $\psi_u$ is a Legendre function
on the interior $I_u$ of its effective domain: $0\in I_u$, $\psi_u$ is differentiable and
strictly convex there, and boundary means in the effective domain of $\psi_u^*$ are limits of
the derivatives $\psi_u'(t)$. Then, for every $r>0$,
\begin{equation}\label{eq:a4-mean-local-dual}
 C_{\mathrm{mean},\mathrm{all},r}(\pi)
 =\sup_{\norm u=1}\ \sup_{\substack{t\in I_u\setminus\{0\}\\0<r_u(t)\le r}}
 \frac{\psi_u'(t)^2}{2r_u(t)},
\end{equation}
where a constant direction contributes zero and endpoint values are understood by closure.
If, in addition, $\E_\pi e^{\tau\norm{\theta}}<\infty$ for some $\tau>0$ and
$\operatorname{Cov}_\pi(\theta)\succ0$, then
\begin{equation}\label{eq:a4-mean-local-cov}
 C_{\mathrm{mean},\mathrm{all},r}(\pi)
 =\lmax\!\big(\operatorname{Cov}_\pi(\theta)\big)+O(\sqrt r)
 \qquad(r\downarrow0).
\end{equation}
For $\pi=N(\mu,\Sigma)$ the equality is exact for every $r>0$:
$C_{\mathrm{mean},\mathrm{all},r}=\lmax(\Sigma)$.
\end{proposition}

\begin{proof}
Fix $u$ and write $I_u(y)=\psi_u^*(y)$. The Gibbs variational principle gives
$\KL(q\|\pi)\ge I_u(\E_qX_u)$. Conversely, the exponential tilt
\[
 \frac{\dd q_{u,t}}{\dd\pi}=e^{tX_u-\psi_u(t)}
\]
has mean $\E_{q_{u,t}}X_u=\psi_u'(t)$ and entropy
$\KL(q_{u,t}\|\pi)=r_u(t)=I_u(\psi_u'(t))$. The Legendre hypothesis and endpoint
approximation therefore identify the entropy infimum for every accessible projected mean. If a
tilt is not itself in $\mathcal P_1(\R^d)$, condition it on increasing Euclidean balls. Since
$t$ is interior, the projected means and entropies of these $\mathcal P_1$ laws converge to
those of the tilt; first use strict entropy slack and then take closure at the boundary.
Taking first the supremum over projected means with entropy at most $r$, and then over $u$, proves
\eqref{eq:a4-mean-local-dual}.

Under the exponential-moment hypothesis, the directional cumulants are uniform for $\norm u=1$
and $t$ in a common neighbourhood of zero. Uniformly in $u$,
\[
 \psi_u(t)=\tfrac12v_ut^2+O(\abs t^3),\quad
 \psi_u'(t)=v_ut+O(t^2),\quad
 r_u(t)=\tfrac12v_ut^2+O(\abs t^3),
 \qquad v_u=u^\top\operatorname{Cov}_\pi(\theta)u.
\]
Positive definiteness and monotonicity of $r_u$ away from zero force
$\abs t=O(\sqrt r)$ on a sufficiently small entropy sublevel. Hence the ratio in
\eqref{eq:a4-mean-local-dual} is $v_u+O(\sqrt r)$ uniformly; maximizing over $u$ proves
\eqref{eq:a4-mean-local-cov}. For a Gaussian,
$\psi_u(t)=\tfrac12t^2u^\top\Sigma u$ and $r_u(t)=\tfrac12t^2u^\top\Sigma u$, so the ratio is
$u^\top\Sigma u$ for every admissible nonzero $t$.
\end{proof}

This proposition is independently agent-certified in
\path{solutions/prop-a4-mean-local.tex}. In particular, the local formula does not contradict
the prior-scale global logistic result below:
small entropy forces $t$ to zero and sees posterior covariance, whereas the unrestricted global
supremum may be attained only along remote tilts.

The clean hierarchy is
\begin{equation}\label{eq:a4-hierarchy}
  C_{\mathrm{mean},\mathcal Q}(\pi)\ \le\ C_{\mathcal Q}(\pi)\ \le\ \CTCI(\pi),
\end{equation}
and the inequalities can be very loose. A posterior can fail global $T_2$ (and have
$C_{\mathcal Q}=\infty$, Warning~\ref{warn:a4-heavytail}) yet have a finite, small mean constant
for the specific linear functionals the statistician reports. For many VB diagnostics
$C_{\mathrm{mean},\mathcal Q}$ is therefore the \emph{right primary object}, with the $\Wtwo$
constant a stronger but frequently unattainable companion; the same construction gives a
per-functional constant $C_{f,\mathcal Q}$ for any scalar observable $f(\theta)$.

\begin{proposition}[Global prior-scale rigidity for Gaussian variational families]
\label{prop:a4-logistic-global}
Let $\pi$ be the Gaussian-prior binary-logistic posterior of
Proposition~\ref{prop:a2-logistic-global}. If $\mathcal Q$ contains all translates
$N(m,S)$ of one fixed $S\succ0$, then
\[
 C_{\mathrm{mean},\mathcal Q}=C_{\mathcal Q}
 =C_{\mathrm{mean},\mathrm{all}}=\CTCI=\lmax(\Sigmaz).
\]
\end{proposition}

\begin{proof}
The upper bounds follow from $C_{\mathrm{mean},\mathcal Q}\le C_{\mathcal Q}\le\CTCI$ and
Proposition~\ref{prop:a2-logistic-global}. For $q_{tv}=N(tv,S)$, the quadratic prior and the
at-most-linear likelihood give
\[
 2\KL(q_{tv}\|\pi)=t^2v^\top\Sigmaz^{-1}v+O(\abs t),
\]
whereas both $W_2^2(q_{tv},\pi)$ and the squared mean displacement are at least
$t^2\norm v^2+O(\abs t)$. Optimizing over $v$ supplies the reverse bound
$\lmax(\Sigmaz)$.
\end{proof}

\noindent The independently agent-certified standalone proof, including the full Gaussian
translation KL expansion and both equality chains, is
\path{solutions/prop-a4-logistic-global.tex}.

The mean/transport gap can also be strict: for the centred scale family
$\mathcal Q=\{N(0,s^2):s>0\}$ and symmetric target $\pi\propto e^{-\abs x^p}$, $1\le p<2$,
$C_{\mathrm{mean},\mathcal Q}=0$ while $C_{\mathcal Q}=\infty$.

\subsection{Three local VI constants: fit, misspecification, and optimization error}
\label{subsec:a4-local-three}

The phrase ``local VI constant'' hides three different ratios. The next three propositions make
their hypotheses and baselines explicit. Let $q_\eta$, $\eta\in\R^p$, be an identifiable
$\mathcal P_2$ family, put $K(\eta)=\KL(q_\eta\|\pi)$,
$D(\eta)=W_2^2(q_\eta,\pi)$, and
$M(\eta)=\norm{\E_{q_\eta}\theta-\E_\pi\theta}^2$. A ``locally isolated sublevel'' means that
for some $\rho_0>0$ the entire family sublevel under discussion is contained in one compact
parameter chart on which $K$ is continuous, the stated optimizer is unique, and equivalently
every neighborhood $V$ of that optimizer satisfies
$\inf_{\eta\notin V,\ K(\eta)\le\delta_{\mathcal Q}+\rho_0}K(\eta)>\delta_{\mathcal Q}$.
This is the finite-dimensional coercivity/separation clause that rules out a second remote
competitor with asymptotically the same KL value.

\begin{proposition}[Well-specified local constant]
\label{prop:a4-local-wellspecified}
Suppose $q_0=\pi$, the sublevel $\{K\le\rho_0\}$ is locally isolated at $0$, and the following
expansions hold uniformly as $h\to0$:
\[
 2K(h)=h^\top Fh+O(\norm h^3),\quad
 D(h)=h^\top Gh+O(\norm h^3),\quad
 \E_{q_h}\theta-\E_\pi\theta=Bh+O(\norm h^2),                 \tag{A4.1}
\]
with $F\succ0$. Then
\[
 C_{\mathcal Q,\rho}(\pi)
 =\lmax(F^{-1/2}GF^{-1/2})+O(\sqrt\rho),
\]
\[
 C_{\mathrm{mean},\mathcal Q,\rho}(\pi)
 =\lmax(F^{-1/2}B^\top BF^{-1/2})+O(\sqrt\rho)
 \qquad(\rho\downarrow0).
\]
\end{proposition}

\begin{proof}
Positive definiteness of $F$, the first expansion, and sublevel isolation imply
$\norm h=O(\sqrt\rho)$ whenever $K(h)\le\rho$. Dividing the quadratic numerator expansions by
$h^\top Fh+O(\norm h^3)$ gives the corresponding generalized Rayleigh quotient with a uniform
$O(\sqrt\rho)$ error. The upper bounds follow by maximizing on the sublevel; the lower bounds
follow by taking arbitrarily small multiples of a top generalized eigenvector. The mean statement
uses $M(h)=h^\top B^\top Bh+O(\norm h^3)$.
\end{proof}

\begin{remark}[A sufficient Wasserstein tangent construction --- unreviewed addition]
\label{rem:a4-poisson-tangent}
One standard route to the quadratic coefficient $G$, beyond the independently reviewed
proposition, is to assume that the scores
$s_i=\partial_{\eta_i}\log q_\eta|_0$ are centred elements of
$L^2(\pi)\cap H^{-1}(\pi)$, that the weak Poisson equations
$-\nabla\!\cdot(\pi\nabla\phi_i)=s_i\pi$ admit finite-energy solutions, and that the required
Wasserstein differentiability and uniform remainder estimates hold. Then
$F_{ij}=\int s_is_j\,\dd\pi$ and
$G_{ij}=\int\langle\nabla\phi_i,\nabla\phi_j\rangle\,\dd\pi$ are the Fisher and Wasserstein
tangent Gram matrices. This identification is useful, but it requires its own regularity proof and
is not part of the reviewed statement of Proposition~\ref{prop:a4-local-wellspecified}.
\end{remark}

\begin{proposition}[Misspecified raw sublevel]
\label{prop:a4-local-misspecified}
Suppose $K$ has a unique interior minimizer $\eta_\star$, let
$\delta=K(\eta_\star)>0$, and assume the sublevel $\{K\le\delta+\rho_0\}$ is locally isolated.
For $h=\eta-\eta_\star$, suppose uniformly
\[
 2K(\eta_\star+h)=2\delta+h^\top F_\star h+O(\norm h^3),
 \quad F_\star\succ0,
\]
\[
 D(\eta_\star+h)=D_\star+b_D^\top h+O(\norm h^2),\qquad
 M(\eta_\star+h)=M_\star+b_M^\top h+O(\norm h^2).
\]
Then
\begin{align}
 C_{\mathcal Q,\delta+\rho}(\pi)
 &=\frac{D_\star}{2\delta}
   +\frac{\sqrt{2\rho}}{2\delta}\norm{F_\star^{-1/2}b_D}+O(\rho),\label{eq:a4-misspec-W2}\\
 C_{\mathrm{mean},\mathcal Q,\delta+\rho}(\pi)
 &=\frac{M_\star}{2\delta}
   +\frac{\sqrt{2\rho}}{2\delta}\norm{F_\star^{-1/2}b_M}+O(\rho).
   \label{eq:a4-misspec-mean}
\end{align}
If the relevant linear term vanishes, the correction is only $O(\rho)$. A boundary optimizer or
multiple KL minimizers requires a different tangent cone or a maximum over optimizer charts and
is not covered by this statement.
\end{proposition}

\begin{proof}
The KL sublevel is, up to an $O(\rho)$ radial perturbation, the ellipsoid
$h^\top F_\star h\le2\rho$. Expanding the denominator about $2\delta$ shows that its first
variation is quadratic, whereas the numerator has the displayed linear variation. Maximizing
$b^\top h$ on that ellipsoid gives
$\sqrt{2\rho}\norm{F_\star^{-1/2}b}$; all numerator quadratic terms, denominator variation, and
the boundary perturbation contribute $O(\rho)$. The optimizer itself supplies the baseline.
\end{proof}

The tangent constant becomes meaningful again if numerator and denominator are both centred at
the variational optimizer. Define, for the unique optimizer $q_\star=q_{\eta_\star}$,
\begin{align*}
 C^{\mathrm{ex}}_{\mathcal Q,\rho}(q_\star)
 &:=\sup_{0<K(\eta)-\delta\le\rho}
 \frac{W_2^2(q_\eta,q_\star)}{2\{K(\eta)-\delta\}},\\
 C^{\mathrm{ex}}_{\mathrm{mean},\mathcal Q,\rho}(q_\star)
 &:=\sup_{0<K(\eta)-\delta\le\rho}
 \frac{\norm{\E_{q_\eta}\theta-\E_{q_\star}\theta}^2}
 {2\{K(\eta)-\delta\}}.
\end{align*}

\begin{proposition}[Optimizer-centred excess-KL constant]
\label{prop:a4-local-excess}
Under the locally isolated interior-minimizer and KL hypotheses of
Proposition~\ref{prop:a4-local-misspecified}, suppose additionally
\[
 W_2^2(q_{\eta_\star+h},q_\star)=h^\top G_\star h+O(\norm h^3),\qquad
 \E_{q_{\eta_\star+h}}\theta-\E_{q_\star}\theta=B_\star h+O(\norm h^2).
\]
Then
\[
 C^{\mathrm{ex}}_{\mathcal Q,\rho}(q_\star)
 =\lmax(F_\star^{-1/2}G_\star F_\star^{-1/2})+O(\sqrt\rho),
\]
\[
 C^{\mathrm{ex}}_{\mathrm{mean},\mathcal Q,\rho}(q_\star)
 =\lmax(F_\star^{-1/2}B_\star^\top B_\star F_\star^{-1/2})+O(\sqrt\rho).
\]
These are optimization-error constants around $q_\star$, not approximation-error certificates
for $q_\star$ versus $\pi$; the two contributions can only be combined, for example, through
$W_2(q_\eta,\pi)\le W_2(q_\eta,q_\star)+W_2(q_\star,\pi)$.
\end{proposition}

\begin{proof}
Subtracting the optimizer value removes the entropy baseline:
$2\{K(\eta_\star+h)-\delta\}=h^\top F_\star h+O(\norm h^3)$. The same generalized-Rayleigh
argument as in Proposition~\ref{prop:a4-local-wellspecified} applies.
\end{proof}

\begin{example}[Exact Gaussian calibration of all three regimes]
\label{ex:a4-gaussian-local}
Let $\pi=N(\mu,\Sigma)$, $\Sigma\succ0$, and
$\mathcal Q_S=\{q_m=N(\mu+m,S):m\in\R^d\}$ for fixed $S\succ0$. Put
\[
 \delta_S=\frac12\!\left\{\operatorname{tr}(\Sigma^{-1}S)-d
       +\log\frac{\det\Sigma}{\det S}\right\},
\quad
 D_S=\operatorname{tr}\!\left[S+\Sigma
       -2(\Sigma^{1/2}S\Sigma^{1/2})^{1/2}\right],
\quad \lambda=\lmax(\Sigma).
\]
Then
\[
 2\KL(q_m\|\pi)=2\delta_S+m^\top\Sigma^{-1}m,\qquad
 W_2^2(q_m,\pi)=D_S+\norm m^2.
\]
If $S\ne\Sigma$, then for every $\rho\ge0$,
\begin{equation}\label{eq:a4-gaussian-misspec}
 C_{\mathcal Q_S,\delta_S+\rho}
 =\max\!\left\{\frac{D_S}{2\delta_S},
 \frac{D_S+2\rho\lambda}{2(\delta_S+\rho)}\right\},
 \qquad
 C_{\mathrm{mean},\mathcal Q_S,\delta_S+\rho}
 =\lambda\frac{\rho}{\delta_S+\rho}.
\end{equation}
For every $\rho>0$ the optimizer-centred transport and mean excess constants are both exactly
$\lambda$. If $S=\Sigma$, the family is well specified and both ordinary localized constants
are exactly $\lambda$ for every $\rho>0$.
\end{example}

\begin{proof}
Write $s=m^\top\Sigma^{-1}m$. At fixed $s$, the largest possible $\norm m^2$ is $\lambda s$,
and the KL sublevel is $0\le s\le2\rho$. The raw transport ratio is therefore the maximum of
$(D_S+\lambda s)/(2\delta_S+s)$ on an interval; its derivative has constant sign, so a maximum
occurs at an endpoint. The mean ratio $\lambda s/(2\delta_S+s)$ is increasing. Centring at
$q_0=N(\mu,S)$ removes $D_S$ and $\delta_S$, leaving
$\norm m^2/(m^\top\Sigma^{-1}m)$, whose supremum is $\lambda$.
\end{proof}

The three propositions and Example~\ref{ex:a4-gaussian-local} are independently agent-certified
in four standalone dossiers:
\begin{center}\small
\path{solutions/prop-a4-local-wellspecified.tex}\qquad
\path{solutions/prop-a4-local-misspecified.tex}\\
\path{solutions/prop-a4-local-excess.tex}\qquad
\path{solutions/ex-a4-gaussian-local.tex}.
\end{center}
They replace the ambiguous use of one tangent formula for three inequivalent questions.

\subsection{Front 3a: heavy tails need modified transport, not ordinary \texorpdfstring{$T_2$}{T2}}
\label{subsec:a4-modified}

Ordinary $T_2$ is Gaussian-tail technology: by Theorem~\ref{thm:heavy-tail-no-lsi} a target with
merely exponential, sub-exponential, or polynomial tails satisfies no LSI and no $T_2$, so
$\Wtwo^2\lesssim\KL$ \emph{cannot} hold uniformly --- this is structural, not technical, and is
the global reason Warning~\ref{warn:a4-heavytail} happens. The right replacement, the $\Wtwo$
analogue of A3, is a \emph{modified} transport inequality with a cost $\alpha$ that is quadratic
near $0$ but grows more slowly at infinity:
\begin{equation}\label{eq:a4-modified}
  \mathcal T_\alpha(q,\pi)\ \le\ C_\alpha\,\KL(q\,\|\,\pi)
  \qquad\text{on}\ \mathcal Q\ \text{or on KL sublevels},
\end{equation}
where for $\pi\propto e^{-\abs x^p}$ ($1<p<2$) one expects $\alpha(r)\sim r^2$ near $0$ and
$\alpha(r)\sim r^{p}$ at infinity, sliding to linear cost for exponential tails and to weighted
$W_1$-type controls for polynomial tails where even $T_1$ fails globally. The modified-transport
and modified-log-Sobolev toolbox is exactly the one developed for these tail classes
\cite{BobkovLedoux1997,GentilGuillinMiclo2005,GozlanLeonard2007}, together with the weighted
Csisz\'ar--Kullback--Pinsker inequalities that tie tail integrability of $\pi$ to weighted
total-variation and Wasserstein bounds \cite{BolleyVillani2005}. The A4 deliverable is to turn
the Tier-4 negative verdict (Family~\ref{fam:heavy-tail-classification}) into a positive,
posterior-specific \emph{classification}: for each robust-Bayes prior (Laplace, Student-$t$,
horseshoe, scale mixtures) identify the strongest cost $\alpha$ for which \eqref{eq:a4-modified}
holds, on $\mathcal Q$ or on KL sublevels.

The one-dimensional log-concave case already has an exact published classification.  For
$1\le p\le2$, put
\[
 \theta_p(t)=
 \begin{cases}
 t^2,&\abs t\le1,\\[1mm]
 \dfrac{2}{p}\abs t^p+1-\dfrac{2}{p},&\abs t\ge1.
 \end{cases}
\]

\begin{theorem}[One-dimensional modified-transport classification]
\label{thm:a4-modified-transport-1d}
Let $\mu$ be a nonatomic, full-support log-concave probability law on $\R$.  Let
$\alpha:\R\to[0,\infty)$ be even, continuous, nondecreasing and superadditive on $\R_+$, with
$\alpha(0)=0$ and $\alpha(t)=t^2$ for $\abs t\le1$, and assume additionally that $\alpha$ is
convex.  A scaled global inequality
\[
 \mathcal T_{\alpha(a\,\cdot)}(\nu,\mu)\le \KL(\nu\|\mu)
 \qquad\text{for every probability law }\nu
\]
holds for some $a>0$ if and only if
$\int\exp(\alpha(bx))\,\mu(\dd x)<\infty$ for some $b>0$.
Consequently, if $\pi_p(\dd x)\propto e^{-\abs x^p}\dd x$ with $1\le p\le2$, then
\[
 \mathcal T_{\theta_p(a_p\,\cdot)}(\nu,\pi_p)
 \le \KL(\nu\|\pi_p)
\]
for some $a_p>0$.  In contrast, a polynomial-tail law admits no global
transport--entropy inequality with a nontrivial unbounded convex cost.
\end{theorem}

\noindent This is Gozlan's real-line criterion \cite[Theorem~57 and
Propositions~16, 26]{Gozlan2007RealLine}.  The positive conclusion is one-dimensional and
existential: it gives neither an optimal $a_p$ nor a dependent, likelihood-tilted, or
family-restricted posterior theorem.  The negative conclusion explains why Student-$t$ and
horseshoe targets require a weighted, weak, restricted, or KL-sublevel formulation rather than
a global convex-cost substitute for $T_2$.

\subsection{Front 3b: multimodality --- the lower bound is family-dependent}
\label{subsec:a4-multimodal}

The Tier-5 intuition is correct --- multimodality can make functional-inequality constants enormous
(Theorem~\ref{thm:separated-mixture-lower}) --- but every A4 lower bound must be stated relative to
$\mathcal Q$. To keep the geometry and the numerical target identical, let
\[
 \pi_a=\tfrac12N(-a,\sigma^2)+\tfrac12N(a,\sigma^2),
 \qquad D:=2a
\]
have centres $\pm a$ and separation $D$.

\emph{Component-reweighting candidate.} The law
$q_\eps=(\tfrac12+\eps)N(-a,\sigma^2)+(\tfrac12-\eps)N(a,\sigma^2)$ is an admissible lower witness
only for a family that contains such mixtures. Data processing from the latent component label
gives
\[
 \KL(q_\eps\|\pi_a)\le
 \KL\big((\tfrac12+\eps,\tfrac12-\eps)\|(\tfrac12,\tfrac12)\big)
 =2\eps^2+O(\eps^4).
\]
At fixed $a/\sigma$, $W_2^2(q_\eps,\pi_a)$ is quadratic as $\eps\to0$; a roughly
$\eps D^2$ mass-transfer regime is expected only beyond an overlap-dependent crossover. If that
crossover is exponentially small in $a^2/\sigma^2$, the resulting ratios provide the familiar
exponential candidate mechanism. This paragraph is a heuristic and a proposed lower-witness
calculation, not a proof that the witness is worst or that $\CTCI$ has a specified prefactor.

\emph{Gaussian mode-collapse lower witness (analytic draft).} A single Gaussian cannot represent the preceding
reweighting. Take instead $q_a=N(a,\sigma^2)$. Then
$\KL(q_a\|\pi_a)\to\log2$, and a one-dimensional quantile calculation gives
$W_2^2(q_a,\pi_a)=D^2/2+O(D\sigma+\sigma^2)$ as $a/\sigma\to\infty$. Consequently any Gaussian
family containing $q_a$ obeys the one-sided lower bound
\begin{equation}\label{eq:a4-modecollapse}
  C_{\mathcal Q}(\pi_a)\ \ge\ \frac{D^2}{4\log2}\,(1+o(1)).
\end{equation}
This proves only quadratic growth of a lower witness. It gives no matching upper bound on
$C_{\mathcal Q}$ and therefore does not show that the restricted constant is merely polynomial.
Likewise, a richer family containing $q_\eps$ inherits its computed lower ratios, but their sharp
separation asymptotic remains open. The three constants $\CTCI(\pi_a)$,
$C_{\mathcal Q}(\pi_a)$, and the localized $C_{\mathcal Q,r}(\pi_a)$ must therefore be audited
separately rather than assigned barrier scales by analogy.

\begin{remark}[Symmetry quotienting is not optional]
\label{rem:a4-quotient}
When the modes are $K!$ relabelings (mixture label switching), a raw $\Wtwo$ between parameter
vectors measures coordinate non-identifiability, not inferential error: $q$ and $\pi$ can differ
only by a permutation yet be reported as far apart. The correct transport problem is then on the
quotient $\Theta/S_K$, $W_{2,\Theta/S_K}^2(q,\pi)\le2C_{\mathrm{quot}}\KL(q\,\|\,\pi)$, or on
permutation-invariant observables (the induced mixture density). This is exactly the
symmetry-quotient escape of Remark~\ref{rem:label-switching}, made quantitative by A5
(Section~\ref{sec:a5}, Theorem~\ref{thm:a5-monotone}): A4 should separate \emph{physical}
multimodality (genuine $C_{\mathcal Q}$ blow-up that survives the quotient) from
\emph{symmetry-induced} multimodality (removed by it).
\end{remark}

\begin{lemma}[Symmetrizing a variational law]
\label{lem:a4-symmetrization}
Let a finite group $G$ act isometrically on $\Theta$, let $\pi$ be $G$-invariant, and for
$q\ll\pi$ in $\mathcal P_2(\Theta)$ put
$\bar q=|G|^{-1}\sum_{g\in G}g_\#q$. Then every integrable $G$-invariant observable $f$ obeys
\[
 \E_{\bar q}f=\E_qf,
\]
while
\begin{equation}\label{eq:a4-symmetrization}
 \KL(\bar q\|\pi)\le\KL(q\|\pi),
 \qquad W_2^2(\bar q,\pi)\le W_2^2(q,\pi).
\end{equation}
The conclusion applies inside a variational optimization only when $\mathcal Q$ is closed under
this averaging (or when the symmetrized family is admitted); it does not by itself order the
ratio of the two decreasing quantities in \eqref{eq:a4-symmetrization}.
\end{lemma}

\begin{proof}
Invariant observables are unchanged by every pushforward $g_\#q$. Convexity of relative entropy
in its first argument and invariance of $\pi$ give
$\KL(\bar q\|\pi)\le |G|^{-1}\sum_g\KL(g_\#q\|\pi)=\KL(q\|\pi)$. For each $g$, push an optimal
coupling between $q$ and $\pi$ through the isometry $(g,g)$; it couples $g_\#q$ to $\pi$ with the
same cost. Averaging these couplings couples $\bar q$ to $\pi$, proving the Wasserstein bound.
\end{proof}

\noindent This lemma is independently agent-certified in
\path{solutions/lem-a4-symmetrization.tex}.

\subsection{The posterior-scale certificate still to prove}
\label{subsec:a4-theorem}

The honest target is not ``$T_2$ certifies VB'' --- that is immediate from \eqref{eq:a4-bridge}
--- but a \emph{restricted, localized, observable-aware} certificate at posterior scale.

\begin{question}[When do posterior-scale VB error bars hold?]
\label{q:a4-certificate}
For the Gaussian-prior convex GLM posterior of Family~\ref{fam:glm-posterior} and a specified
mean-field or Gaussian family $\mathcal Q$, put
$A_{\bar W}:=\Sigmaz^{-1}+X^\top\bar W X$ and suppose the KL minimizer $q^\star$ exists. Find
explicit, verifiable conditions under which there are $\rho_0>0$ and $K<\infty$ such that, for
every $0<\rho\le\rho_0$,
\[
  C_{\mathcal Q,\delta_{\mathcal Q}+\rho}(\pi)
  \ \le\ K\lmax(A_{\bar W}^{-1})+R_2(\rho),
  \qquad
  C_{\mathrm{mean},\mathcal Q,\delta_{\mathcal Q}+\rho}(\pi)
  \ \le\ K\lmax(A_{\bar W}^{-1})+R_{\rm mean}(\rho).
\]
At minimum, state and verify a coercivity--tail contract: compactness of the relevant family
sublevel, uniform integrability in $\mathcal P_2$, bulk Hessian control at the scale
$A_{\bar W}$, and quantitative coupling and mean-tail bounds that define the finite remainders
$R_2$ and $R_{\rm mean}$. The remainders must include the misspecification baselines and possible linear terms of
Proposition~\ref{prop:a4-local-misspecified}; they vanish at the appropriate rate only under an additional
well-specified/local-centering hypothesis. These bounds yield the
certificates $\Wtwo^2(q^\star,\pi)\le2C_{\mathcal Q,r}\,\KL(q^\star\,\|\,\pi)$ and
$\norm{\E_{q^\star}\theta-\E_\pi\theta}^2\le2C_{\mathrm{mean},\mathcal Q,r}\,\KL(q^\star\,\|\,\pi)$,
provided a separate computable upper bound on $\KL(q^\star\,\|\,\pi)$ is supplied
(Remark~\ref{rem:a4-elbo}).
The global constants are not part of this claim: under the location-rich hypothesis of
Proposition~\ref{prop:a4-logistic-global}, both are exactly $\lmax(\Sigmaz)$.
\end{question}

\noindent The proof route mirrors A1/A2: Brascamp--Lieb (Theorem~\ref{thm:brascamp-lieb}) or a
Lyapunov/local-curvature argument supplies the bulk curvature
$\Sigmaz^{-1}+X^\top\bar W X$ that the global floor of Theorem~\ref{thm:glm-fi} discards;
the stated compactness and uniform-$\mathcal P_2$ clauses prevent escape inside the family, while
the tail remainder records rather than suppresses the part outside the bulk
(Warning~\ref{warn:a4-heavytail}); and the mean constant is calibrated by the exact
Donsker--Varadhan dual of Subsection~\ref{subsec:a4-mean}. Lower witnesses must belong to the
chosen family; covariance alone is not a lower bound for a restricted mean constant.

\subsection{Sub-targets and what remains open}
\label{subsec:a4-open}

The program splits into five research questions, ordered by difficulty.

\begin{question}[Restricted and localized GLM transport constants]
\label{q:a4-restricted}
Compute or sharply upper-bound $C_{\mathcal Q,\delta_{\mathcal Q}+\rho}(\pi)$ for the GLM
posteriors of Section~\ref{sec:tier3} (and for sparse-regression and hierarchical posteriors),
proving posterior-scale constants $\asymp\lmax[(\Sigmaz^{-1}+X^\top\bar W X)^{-1}]$ in the bulk
where the global $C_{\mathcal Q}=\CTCI=\lmax(\Sigmaz)$ is prior-scale for location-rich Gaussian
families. Use the tangent generalized eigenvalue of
Proposition~\ref{prop:a4-local-wellspecified} as the
well-specified small-$\rho$ calibration and state the coercivity/tail assumption that makes the
sublevel finite. The exact localized supremum
\eqref{eq:a4-cqr} is not a standard output of either transport theory or VI theory; this is the
$\Wtwo$ analogue of the flagship A1 (Question~\ref{q:glm-data-informed}) and the entry point.
\end{question}

\begin{question}[First-moment constants]
\label{q:a4-mean}
Use the exact global and localized cumulant duals \eqref{eq:a4-mean-dual} and
\eqref{eq:a4-mean-local-dual} for unrestricted means and compute localized
or per-functional constants for variational families. Explain when the family can realize the
worst entropy tilt, and classify gaps such as the centred scale-family example where
$C_{\mathrm{mean},\mathcal Q}=0$ but $C_{\mathcal Q}=\infty$. For logistic regression the global
location-rich constant is already fixed by Proposition~\ref{prop:a4-logistic-global}, so the
posterior-scale question must be localized.
\end{question}

\begin{question}[Modified transport for heavy-tailed posteriors]
\label{q:a4-modified}
The one-dimensional generalized-normal branch is solved by
Theorem~\ref{thm:a4-modified-transport-1d}, while a polynomial-tail Student or horseshoe law
admits no global transport--entropy inequality with a nontrivial unbounded convex cost.  For a
first residual target, take
\[
 \pi_{\nu,n}(\dd x)\propto
 \left(1+\frac{x^2}{\nu}\right)^{-(\nu+1)/2}
 \operatorname{sigmoid}(x)^n\dd x,
 \qquad \nu>2,\quad n\in\mathbb N,\ n\ge1,
\]
and the fixed-scale Gaussian location family
$\mathcal Q_s=\{q_m=N(m,s^2):m\in\R\}$, $s>0$.  Writing
$K(m)=\KL(q_m\|\pi_{\nu,n})$, $\delta=\min_mK(m)$, and
$M_\rho=\{m:K(m)\le\delta+\rho\}$, first prove continuity and reverse-KL coercivity of $K$,
$\delta>0$, and nonempty compactness of every complete sublevel $M_\rho$.  Then obtain explicit
matching upper and lower bounds, including
the sharp large-$\rho$ order, for
\[
 C_{\nu,n,s}(\rho)
 =\sup_{m\in M_\rho}
 \frac{W_2^2(q_m,\pi_{\nu,n})}{2K(m)},
 \qquad \rho\ge0.
\]
Finiteness must come from this named family's compact sublevels and the $\mathcal P_2$ tail
contract $\nu>2$, not from a KL cutoff alone.  A horseshoe target must instead specify a
weighted, weak, bounded, or otherwise finite cost together with its family-coercivity contract.
\end{question}

\begin{question}[Family-aware multimodal lower bounds]
\label{q:a4-multimodal}
For the Gaussian mixture with centres $\pm a$ and separation $D=2a$, certify the component-
reweighting and mode-collapse lower witnesses of Subsection~\ref{subsec:a4-multimodal}. Determine,
rather than assume, their crossover and sharp separation dependence; in particular, seek matching
upper bounds before asserting an asymptotic for $C_{\mathcal Q}$. Pair symmetry-induced cases with
the quotient escape of Remark~\ref{rem:a4-quotient}, separating them from physical multimodality
(A5, Section~\ref{sec:a5}).
\end{question}

\begin{question}[Coupling the constant to a computable KL certificate]
\label{q:a4-elbo}
Combine $C_{\mathcal Q,r}$ with a computable upper bound on $\KL(q^\star\,\|\,\pi)$
(equivalently an \emph{upper} bound on $\log Z$ combined with the ELBO) to produce a genuinely numerical,
end-to-end VB error bar (Remark~\ref{rem:a4-elbo}). This is what turns
Question~\ref{q:a4-certificate} from an inequality between constants into a deliverable a practitioner
can report, and is the bridge from the functional-inequality theory of this Part to the
algorithmic VI literature \cite{LambertEtAl2022,ArneseLacker2024,JiangChewiPooladian2023}.
\end{question}

\begin{remark}[Relation to the rest of the agenda]
\label{rem:a4-relations}
A4 is the transport/variational face of the agenda, and it \emph{consumes} the other targets more
than it competes with them. Its bulk constant is the $\Wtwo$ analogue of the finite-sample flagship
A1 (Section~\ref{sec:a1}, Conjecture~\ref{conj:a1}), at the same inverse-observed-information scale
$\lmax[(\Sigmaz^{-1}+X^\top\bar W X)^{-1}]$; in the large-sample limit that scale is the
Fisher-information constant of A2 (Section~\ref{sec:a2}, Conjecture~\ref{conj:a2},
Equation~\eqref{eq:a2-target}), so a sharp localized $C_{\mathcal Q,r}$ should reproduce
$\tfrac1n\lmax(I(\theta_0)^{-1})$. The heavy-tailed obstruction that makes $C_{\mathcal Q}=\infty$
(Warning~\ref{warn:a4-heavytail}) is the transport shadow of A3 (Section~\ref{sec:a3}), and the
modified cost of Question~\ref{q:a4-modified} (Equation~\eqref{eq:a4-modified}) is the $\Wtwo$
counterpart of A3's weighted Poincar\'e inequality. Finally, the family-dependent multimodal lower
bounds of Subsection~\ref{subsec:a4-multimodal} and the quotient escape of
Remark~\ref{rem:a4-quotient} are the variational reading of A5 (Section~\ref{sec:a5},
Theorem~\ref{thm:a5-monotone}), separating physical multimodality from the symmetry-induced
multimodality that A5 quotients away.
\end{remark}
\end{document}
